%% 
%% Copyright 2007-2025 Elsevier Ltd
%% 
%% This file is part of the 'Elsarticle Bundle'.
%% ---------------------------------------------
%% 
%% It may be distributed under the conditions of the LaTeX Project Public
%% License, either version 1.3 of this license or (at your option) any
%% later version.  The latest version of this license is in
%%    http://www.latex-project.org/lppl.txt
%% and version 1.3 or later is part of all distributions of LaTeX
%% version 1999/12/01 or later.
%% 
%% The list of all files belonging to the 'Elsarticle Bundle' is
%% given in the file `manifest.txt'.
%% 
%% Template article for Elsevier's document class `elsarticle'
%% with harvard style bibliographic references

\documentclass[preprint,12pt]{elsarticle}
\usepackage{subcaption}
\usepackage{hyperref}
%% Use the option review to obtain double line spacing
%% \documentclass[preprint,review,12pt]{elsarticle}

%% Use the options 1p,twocolumn; 3p; 3p,twocolumn; 5p; or 5p,twocolumn
%% for a journal layout:
%% \documentclass[final,1p,times]{elsarticle}
%% \documentclass[final,1p,times,twocolumn]{elsarticle}
%% \documentclass[final,3p,times]{elsarticle}
%% \documentclass[final,3p,times,twocolumn]{elsarticle}
%% \documentclass[final,5p,times]{elsarticle}
%% \documentclass[final,5p,times,twocolumn]{elsarticle}

%% For including figures, graphicx.sty has been loaded in
%% elsarticle.cls. If you prefer to use the old commands
%% please give \usepackage{epsfig}


%% The amssymb package provides various useful mathematical symbols
\usepackage{amssymb}
%% The amsmath package provides various useful equation environments.
\usepackage{amsmath}
\usepackage{xcolor}
\usepackage{booktabs}
\usepackage{tabularx}
\newcommand{\revblue}[1]{\textcolor{blue}{#1}}
%% The amsthm package provides extended theorem environments
%% \usepackage{amsthm}

%% The lineno packages adds line numbers. Start line numbering with
%% \begin{linenumbers}, end it with \end{linenumbers}. Or switch it on
%% for the whole article with \linenumbers.
%% \usepackage{lineno}

\journal{Robotics and Autonomous Systems}

\begin{document}

\begin{frontmatter}

\title{Multimodal robot with bio-inspired gait arbitration across simulated and real environments} %% Article title

\author{Samuel Federico Carlos-Rodríguez\fnref{depA}}
\author{Diego Alejandro Caro-Vaca\fnref{depA}}
\author{David Caicedo-Samboní\fnref{depA}}
\author{Sebastián Grajales-Pulgarín\fnref{depA}}
\author{Julián Hurtado-López\corref{cor1}}
\ead{jhurtado@uao.edu.co}
\cortext[cor1]{Corresponding author}
\author{David Fernando Ramírez-Moreno}

\affiliation{organization={Faculty of Engineering and Basic Sciences, Universidad Autónoma de Occidente},
            addressline={Cll 25 \# 115-85 Km 2 Vía Cali - Jamundí},
            city={Cali},
            postcode={760030},
            state={Valle},
            country={Colombia}}

\fntext[depA]{Mechatronic Engineering}

%% Abstract
\begin{abstract}
This work presents the design, simulation, and implementation of a bio-inspired quadruped robotic platform equipped with a hybrid wheel-leg locomotion system and controlled through a neural architecture modeled after basal ganglia mechanisms. The platform integrates sensory modules for environmental perception, arbitration modules for behavioral decision-making, and locomotor modules for gait generation, enabling adaptive responses to diverse stimuli and terrain conditions. Simulation and implementation demonstrate the robot's ability to switch autonomously between locomotion modes ---quadrupedal walking, differential rolling, and omnidirectional motion--- when facing obstacles, appetitive stimuli, or aversive conditions. \revblue{The proposed framework demonstrates the robustness and reliability of a bio-inspired neural architecture for multi-stimulus mode arbitration, providing a dependable algorithmic foundation for autonomous systems operating under uncertainty and dynamic environmental conditions. Across simulated and real-robot trials, the controller consistently avoided critical instability and produced terrain-specific, repeatable locomotion signatures, highlighting its potential as a reliable framework for gait selection and laying the groundwork for scaling toward robust real-world infrastructure inspection and operation in complex, unstructured environments.}
\end{abstract}

%%Graphical abstract
\begin{graphicalabstract}
\includegraphics[width=\textwidth]{GRAPHICAL_ABSTRACT}
\end{graphicalabstract}

%%Research highlights
\begin{highlights}
\item Hybrid wheel-leg quadruped enables adaptive locomotion in complex environments.
\item Multimodal sensing (camera, LiDAR, IMU) drives context-aware gait selection.
\item A bio-inspired neural architecture arbitrates behaviors via winner-take-all dynamics.
\item The controller switches adaptively among rolling, quadrupedal, and omnidirectional modes in response to sensory context.
\item Results show consistent behavior selection under changing stimuli and environmental conditions.
\end{highlights}

%% Keywords
\begin{keyword}

Wheel-legged quadruped \sep Hybrid locomotion \sep Gait mode switching \sep Arbitration \sep Basal ganglia model \sep Multimodal sensing

\end{keyword}

\end{frontmatter}

%% Add \usepackage{lineno} before \begin{document} and uncomment 
%% following line to enable line numbers
%% \linenumbers

%% main text
%%
